\section{Описание}

Требуется реализовать алгоритм поиска подстроки в последовательности элементов (слов в тексте) используя алгоритм Кнута-Морриса-Пратта (KMP). Особенностью этой реализации является построение функции отказов через промежуточное вычисление Z-array (массива префиксов из Z-algorithm), а затем конвертация его в так называемый \enquote{strong prefix function} для оптимизированного поиска.

\subsection{Основная идея}

Алгоритм KMP решает задачу поиска паттерна длины $m$ в тексте длины $n$ за время $O(n + m)$ без предварительной обработки всего текста в памяти. Вместо наивного перебора со сдвигом на каждый символ, KMP использует информацию о структуре паттерна, чтобы пропускать заведомо неудачные сдвиги.

\subsection{Z-algorithm и Strong Prefix Function}

Z-array массива $p$ представляет собой массив длин наибольших префиксов, совпадающих с подстроками, начинающимися на каждой позиции:
\[
z[i] = \max\{k : p[0..k-1] = p[i..i+k-1]\}
\]

Strong prefix function преобразует Z-array так, что для каждой позиции $j$ в паттерне хранится максимальная длина суффикса, который совпадает с префиксом и заканчивается именно в позиции $j$.

Построение strong prefix function из Z-array:
\[
sp[i + z[i] - 1] = \max(sp[i + z[i] - 1], z[i]) \quad \text{для всех } i
\]

\pagebreak

\section{Исходный код}

Реализация разделена на модули: \texttt{utils.hpp/utils.cpp} содержат утилиты для построения структур данных и поиска, а \texttt{lab4.cpp} содержит главную программу.

\subsection{Построение Z-array}

Алгоритм построения Z-array работает за линейное время $O(m)$ и использует технику скользящего окна:

\begin{lstlisting}[language=C++]
std::vector<int> build_z(const std::vector<std::string>& p) {
    int m = (int)p.size();
    std::vector<int> z(m, 0);
    for (int i = 1, l = 0, r = 0; i < m; ++i) {
        if (i <= r) {
            z[i] = std::min(r - i + 1, z[i - l]);
        }
        while (i + z[i] < m && p[z[i]] == p[i + z[i]]) {
            z[i]++;
        }
        if (i + z[i] - 1 > r) {
            l = i;
            r = i + z[i] - 1;
        }
    }
    return z;
}
\end{lstlisting}

Переменные $l$ и $r$ отслеживают границы наиболее правого отрезка, для которого известна информация о совпадении с префиксом. Это позволяет избежать избыточного сравнения и достичь линейности.

\subsection{Построение Strong Prefix Function}

\begin{lstlisting}[language=C++]
std::vector<int> build_strong_pi(const std::vector<int>& z) {
    int m = (int)z.size();
    std::vector<int> sp(m, 0);
    
    for (int i = m - 1; i >= 0; --i) {
        if (z[i] > 0) {
            int end_idx = i + z[i] - 1;
            sp[end_idx] = std::max(sp[end_idx], z[i]);
        }
    }
    return sp;
}
\end{lstlisting}

\subsection{Поиск в потоке}

\begin{lstlisting}[language=C++]
void kmp_stream_search(std::istream& in, 
                       const std::vector<std::string>& pattern, 
                       const std::vector<int>& sp) {
    if (pattern.empty()) return;

    const int m = (int)pattern.size();
    std::vector<Position> buffer(m);
    
    std::string line;
    size_t line_num = 1;
    size_t total_words = 0;
    int j = 0;
    
    std::string current_word;
    current_word.reserve(100); 

    while (std::getline(in, line)) {
        size_t word_idx = 1;
        
        for (size_t i = 0; i <= line.length(); ++i) {
            if (i == line.length() || 
                std::isspace((unsigned char)line[i])) {
                if (!current_word.empty()) {
                    to_lower_inplace(current_word);

                    buffer[total_words % m] = {line_num, word_idx};
                    ++total_words;

                    while (j > 0 && pattern[j] != current_word) {
                        j = sp[j - 1];
                    }
                    if (pattern[j] == current_word) {
                        ++j;
                    }

                    if (j == m) {
                        const auto& start_pos = 
                            buffer[(total_words - m) % m];
                        std::cout << start_pos.line << ", " 
                                  << start_pos.word_idx << "\n";
                        j = sp[m - 1];
                    }

                    current_word.clear();
                    word_idx++;
                }
            } else {
                current_word += line[i];
            }
        }
        line_num++;
    }
}
\end{lstlisting}

Ключевая оптимизация: используется циклический буфер размером $m$, чтобы хранить только последние $m$ позиций слов. Это обеспечивает $O(m)$ использования памяти вместо хранения всех позиций.

Основной цикл KMP выполняет следующие операции:
\begin{enumerate}
    \item Читает строку из потока и парсит её на слова
    \item Для каждого слова выполняет стандартный KMP-шаг: backtracking при несовпадении через функцию отказов
    \item При полном совпадении паттерна выводит позицию и переходит к следующему возможному совпадению
\end{enumerate}

Сложность алгоритма: предобработка паттерна за 2m, следовательно ассимптотически $O(n + m)$, где $n$ — общее число букв в словах текста, $m$ — число букв в словах паттерна.

\pagebreak
